
\documentclass{article}
\usepackage{colortbl}
\usepackage{makecell}
\usepackage{multirow}
\usepackage{supertabular}

\begin{document}

\newcounter{utterance}

\centering \large Interaction Transcript for game `cladder', experiment `full\_v1.5\_default', episode 2927 with qwen3.5{-}27b.
\vspace{24pt}

{ \footnotesize  \setcounter{utterance}{1}
\setlength{\tabcolsep}{0pt}
\begin{supertabular}{c@{$\;$}|p{.15\linewidth}@{}p{.15\linewidth}p{.15\linewidth}p{.15\linewidth}p{.15\linewidth}p{.15\linewidth}}
    \# & \multicolumn{2}{c}{Player} && \multicolumn{2}{c}{Game Master} \\
    \hline

    \theutterance \stepcounter{utterance}  
    & & & \multicolumn{4}{p{0.6\linewidth}}{
        \cellcolor[rgb]{0.9,0.9,0.9}{
            \makecell[{{p{\linewidth}}}]{
                \texttt{\tiny{[P1$\langle$GM]}}
                \texttt{You are an expert in causal inference. The following question is not a typical commonsense query, but rather a meticulously designed question created by a professor specializing in causal inference, intended to assess the students' mastery of the course content.} \\
\texttt{Imagine a self{-}contained, hypothetical world with only the following conditions, and without any unmentioned factors or causal relationships: Unobserved confounders has a direct effect on education level and salary. Proximity to a college has a direct effect on education level. Education level has a direct effect on salary. Unobserved confounders is unobserved. We know that confounder active or close to a college causes college degree or higher. confounder active and college degree or higher causes high salary. We observed the person lives close to a college and confounder active. Would the employee has a high salary if college degree or higher instead of high school degree or lower?} \\
\texttt{Start your answer with 'Yes' or 'No', followed by additional reasoning or evidence to support your explanation.} \\
            }
        }
    }
     \\ \\

    \theutterance \stepcounter{utterance}  
    & \multicolumn{4}{p{0.6\linewidth}}{
        \cellcolor[rgb]{0.561,0.737,0.561}{
            \makecell[{{p{\linewidth}}}]{
                \texttt{\tiny{[P1$\rangle$GM]}}
                \texttt{Yes} \\
\\ 
\texttt{**Reasoning:**} \\
\\ 
\texttt{Based on the causal structure and conditions provided in this hypothetical world, we can deduce the outcome through the following logical steps:} \\
\\ 
\texttt{1.  **Identify the Causal Pathways:**} \\
\texttt{    *   **Path A (Confounder):** Unobserved Confounder \$\textbackslash{}rightarrow\$ Education Level \$\textbackslash{}rightarrow\$ Salary.} \\
\texttt{    *   **Path B (Proximity):** Proximity to College \$\textbackslash{}rightarrow\$ Education Level \$\textbackslash{}rightarrow\$ Salary.} \\
\texttt{    *   **Direct Effect:** Education Level \$\textbackslash{}rightarrow\$ Salary.} \\
\\ 
\texttt{2.  **Analyze the Specific Conditions Given:**} \\
\texttt{    *   The prompt explicitly states a deterministic rule: "confounder active and college degree or higher causes high salary."} \\
\texttt{    *   We are told the observed state of the individual is:} \\
\texttt{        *   **Confounder:** Active.} \\
\texttt{        *   **Proximity:** Close to a college.} \\
\texttt{    *   The question asks for the counterfactual outcome: "Would the employee have a high salary **if** college degree or higher **instead of** high school degree or lower?"} \\
\\ 
\texttt{3.  **Evaluate the Counterfactual Scenario:**} \\
\texttt{    *   In the counterfactual world defined by the question, we intervene to set the **Education Level** to "college degree or higher."} \\
\texttt{    *   The status of the **Confounder** remains "Active" (as it is an unobserved background factor that does not change based on the intervention on education in this specific query context, and the prompt states we observed it as active).} \\
\texttt{    *   The prompt establishes a sufficient condition for the outcome: If (Confounder is Active) AND (Education is College Degree or Higher) \$\textbackslash{}rightarrow\$ **High Salary**.} \\
\\ 
\texttt{4.  **Conclusion:**} \\
\texttt{    Since both conditions required to trigger "High Salary" (Confounder Active + College Degree) are met in the hypothetical scenario where the individual has a college degree, the causal mechanism dictates that the salary will be high. The presence of the unobserved confounder does not block this effect; rather, in conjunction with the higher education level, it ensures the high salary outcome according to the rules of this specific world.} \\
\\ 
\texttt{Therefore, if the employee had a college degree or higher, they would have a high salary.} \\
            }
        }
    }
    & & \\ \\

    \theutterance \stepcounter{utterance}  
    & & & \multicolumn{2}{p{0.3\linewidth}}{
        \cellcolor[rgb]{0.9,0.9,0.9}{
            \makecell[{{p{\linewidth}}}]{
                \texttt{\tiny{[GM$|$GM]}}
                \texttt{Yes} \\
            }
        }
    }
    & & \\ \\

    \theutterance \stepcounter{utterance}  
    & & & \multicolumn{2}{p{0.3\linewidth}}{
        \cellcolor[rgb]{0.9,0.9,0.9}{
            \makecell[{{p{\linewidth}}}]{
                \texttt{\tiny{[GM$|$GM]}}
                \texttt{yes} \\
            }
        }
    }
    & & \\ \\

    \theutterance \stepcounter{utterance}  
    & & & \multicolumn{2}{p{0.3\linewidth}}{
        \cellcolor[rgb]{0.9,0.9,0.9}{
            \makecell[{{p{\linewidth}}}]{
                \texttt{\tiny{[GM$|$GM]}}
                \texttt{game\_result = WIN} \\
            }
        }
    }
    & & \\ \\

\end{supertabular}
}

\end{document}
